\begin{abstract}
A model that recognises it is being tested may behave differently once deployed. Several
detectors try to flag this evaluation awareness, and if they measure the same thing,
they should agree. We analysed five detectors on Qwen 3.6 27B and three on Gemma 4 31B with
EvalAwareBench, which adds test-like cues to otherwise unchanged prompts. Every analysed
detector passes a known-groups check: it flags cued prompts more often. However,
\cv{qwen/null/n} random directions in Qwen's activations pass it too after the same tuning as
the activation probe. We then asked whether the detectors agree beyond the cues,
comparing them only among prompts that carry the same cues (a post hoc analysis). There, self-report and third-person reach
\cv{gemma/sr-tp/kappawithinfrac}\% (Gemma) and \cv{qwen/sr-tp/kappawithinfrac}\% (Qwen) of the
highest Cohen's $\kappa$ their positive rates allow. That maximum is only
\cv{gemma/sr-tp/kappawithinmax} on Gemma and \cv{qwen/sr-tp/kappawithinmax} on Qwen, so most of
their disagreement comes from flagging at different rates. In contrast, among these prompts the Qwen activation probe
agrees with none of the three other detectors it can be compared with, even though it passes
the known-groups check. A known-groups check is therefore not enough to show that a detector
measures evaluation awareness.
\end{abstract}
